% 根据作者更新的 数据加工.txt 整理；500 为目标规模，359,355 为事实数。
\begin{samepage}
\paragraph{源证据规范化与事实抽取。}
\label{app:data_preparation}
源数据加工将异质医学图文记录转化为可复用的事实证据，为学习者条件化的监督构造提供内容基础。每条记录保留视觉输入、原始文本与来源关联，报告样本保留同次检查的多幅影像。原生问答与分类标签通过规则转化为事实陈述，报告依据 Findings 与 Impression 分节组织，图像描述中的事实则由 Qwen3.6-27B 抽取。源记录沿任务形式、视觉能力、图像模态与解剖部位四个维度标注，使证据检索与课程指定的学习需求相对应。

质量控制考察记录的结构完整性、医学相关性与事实完备性，无有效事实的记录不参与监督构造。新增记录依据样本标识、影像路径及 RGB 像素内容去重，并检查与验证集影像的重合。最终源池提供359,355个可复用事实；该数量以事实条目计数，与图~\ref{fig:support_composition}中的源记录数量区分。

\par
\end{samepage}

\paragraph{基于任务的验证集构建。}
验证集以医学任务为基本组织单元，依据任务的输入属性与输出要求，对已有医学视觉问答与报告数据实施分层选样。问答任务按语言、答案类型、问题形式与内容类别分层，多项选择任务进一步考虑正确选项与问题长度；报告生成任务按影像数量等输入属性分层。各验证子集的目标规模统一设为500条，并保留样本标识、影像、问题及原始参考标注。由此形成的任务相关验证集合，用于持续估计学习者在不同任务条件下的能力表现及其学习响应。

\paragraph{基于任务的四维标注与能力诊断。}
四维标注依据医学任务的能力目标及其证据属性，表征验证实例所对应的任务条件（表~\ref{tab:feedback_axes}）。任务形式与视觉能力分别刻画输出结构与主要能力目标，图像模态与解剖部位则刻画任务所依托的医学证据。任务形式根据任务要求与答案格式确定；视觉问答实例的其余三维标签由 Qwen3.6-27B 依据影像、问题及参考答案补充，优先保留已有模态标注。胸部 X-ray 报告任务的四维标签依次为报告生成、整体描述、胸部 X-ray 与胸部。经格式与完整性检查后，依据这些任务属性分组汇总验证表现，刻画学习者在不同能力目标下的能力分布，为后续监督构造与课程分配确定任务重点。

\begin{table}[!htbp]
\centering\small
\caption{基于任务的验证集标注体系。任务形式与视觉能力刻画能力目标，图像模态与解剖部位刻画医学证据属性；各维度共同用于任务层面的能力诊断与课程适配。视觉能力维度包含九类。}
\label{tab:feedback_axes}
\setlength{\tabcolsep}{4pt}
\renewcommand{\arraystretch}{1.12}
\begin{tabularx}{\linewidth}{@{}lYY@{}}
\toprule
维度 & 类别或内容 & 标注依据\\
\midrule
任务形式 & 是非问答、开放短答、开放长答、多项选择、报告生成 & 任务要求与答案格式。\\
视觉能力 & 模态识别、器官识别、定位、空间关系、异常检测、属性判断、比较、整体描述、医学知识 & 问题的主要考查目标；报告任务归为整体描述。\\
图像模态 & CT、MRI、X-ray、超声、病理图像、论文图等 & 优先保留已有标注，缺失时由模型补充。\\
解剖部位或组织 & 胸部、腹部、脑、骨盆、四肢、脊柱及明确的器官或组织 & 影像与文本证据；无法确定时保留 unknown / other。\\
\bottomrule
\end{tabularx}
\end{table}
